\documentclass[12pt]{article}
\usepackage[utf8]{inputenc}
\usepackage[T1]{fontenc}
\usepackage{hyperref}
\title{Elymian Language Resources for K-12 Curriculum Development}
\author{Compiled by Claude Sonnet 5.5 (claude-sonnet-5-5) for Andr\'as Kornai}
\date{}
\begin{document}
\maketitle

\section{Introduction}
Elymian (ISO 639-3 xly; Glottocode elym1237) was the language of the Elymians of western Sicily (Segesta, Eryx). It is extinct and undeciphered, and its classification is unsettled. According to the sources found, about 170 short inscribed pottery fragments (mostly from a votive deposit at Segesta), coin legends and proper names survive, written in Greek letters. The most complete text is a vase from Montedoro, tentatively read as a dedication naming Ata son of Tuka. One Oxford blog post argues that some Segesta inscriptions show a Libyan-Berber language and not Elymian. That is a minority view. Almost nothing exists for teaching, and a curriculum cannot use the language.

\section{Primary Educational Resources}
I found no primers, textbooks or school programmes for Elymian.

\section{Linguistic and Reference Materials}
\begin{itemize}
\item Elymian language overview (Wikipedia mirror): \url{https://wikipedia.com/wiki/ISO_639:xly}
\item Elymians (Wikipedia): \url{https://en.wikipedia.org/wiki/Elymians}
\item Elymian overview: \url{https://expiredknowledge.com/elymian/}
\item Oxford CrossReads, ``A new language in the epigraphic landscape of ancient Sicily'' (a dissenting view): \url{https://crossreads.web.ox.ac.uk/node/1992046}
\end{itemize}
The search summary also named S. Marchesini, ``The Elymian language'', in Tribulato (ed.), \emph{Language and Linguistic Contact in Ancient Sicily} (Cambridge University Press, 2012). It is not linked in the results, so no URL is given.

\section{Religious and Community Resources}
The votive inscriptions from Segesta are the only religious material. I found no community organisations.

\section{Digital and Audio Resources}
I found no audio or software resources. There is a tourist-oriented page on Segesta: \url{https://www.seepalermo.com/segesta.htm}

\section{Curriculum Development Context}
A reasonable use is a short unit on the peoples of ancient Sicily, the Elymian sites and the problem of undeciphered languages. Students could compare Greek-letter inscriptions and discuss why the evidence is so thin. The language cannot be taught.
\end{document}